\chapter{{\fontsize{14pt}{21pt}\selectfont\MakeUppercase{Реализация веб-приложения}}}
\label{cha:impl}

\section{Выбор набора технологий}

Практическая часть реализована на Python 3.12. Серверная часть построена на FastAPI, схемы входных и выходных данных описаны Pydantic, доступ к базе данных выполнен через SQLAlchemy, миграции управляются Alembic~\cite{python-docs,fastapi-docs,pydantic-docs,sqlalchemy-docs,alembic-docs}. В качестве базы данных используется PostgreSQL с расширением pgvector, что позволяет хранить обычные сущности приложения и векторные представления фрагментов в одной СУБД~\cite{postgres-docs,pgvector}.

Для подготовки корпуса используются Beautiful Soup и локальные файлы HTML и Markdown~\cite{beautifulsoup-docs}. HTTP-взаимодействие с Groq API выполняется через httpx~\cite{httpx-docs,groq-docs}. Векторные представления в рабочем режиме формируются локальной моделью \texttt{BAAI/bge-m3} через библиотеку sentence-transformers; тестовый хеширующий модуль разрешён только при \texttt{APP\_ENV=test}~\cite{sentence-transformers-docs}. Автоматические проверки реализованы на pytest~\cite{pytest-docs}. Локальное развёртывание описано через Docker Compose.

Выбранный набор технологий соответствует задачам ВКР: FastAPI обеспечивает маршруты API и HTML-страницы, PostgreSQL хранит корпус и историю, pgvector обеспечивает векторный поиск, Alembic фиксирует физическую схему, а pytest позволяет повторяемо проверить критичные правила.

FastAPI выбран как основа серверной части, потому что приложение сочетает программный интерфейс и HTML-страницы. Через маршруты API пользовательский интерфейс получает список версий, отправляет вопрос, читает историю и выполняет диагностические запросы. Через обычные маршруты приложение отдаёт главную страницу и страницу истории. Для практической части ВКР это удобно: серверная часть и пользовательский интерфейс находятся в одном проекте, но обмен между ними идёт через явные HTTP-контракты.

Pydantic используется для описания входных и выходных структур. Для основного запроса это особенно важно, потому что пользователь должен передать три обязательных поля: вопрос, версию PostgreSQL и режим ответа. Схема \texttt{AskRequest} проверяет длину вопроса, числовой формат версии и принадлежность версии списку поддерживаемых значений. Режим ответа ограничен тремя допустимыми вариантами: кратким, подробным и обучающим. Такая проверка выполняется до запуска поиска фрагментов и не допускает обработку запроса с неопределённым режимом или неподдерживаемой версией.

SQLAlchemy используется как слой работы с реляционной моделью, а Alembic -- для управляемого изменения физической схемы. В проекте есть связанные сущности: версии, документы, фрагменты, векторные представления, история и источники запросов. ORM-модели описывают эти связи в коде, а миграции фиксируют структуру таблиц, ограничения допустимых значений, поле \texttt{vector(1024)}, JSONB для обучающего результата и индекс для векторного поиска. Поэтому схема базы данных воспроизводится при развёртывании, а не создаётся вручную.

PostgreSQL с расширением pgvector выбран потому, что приложению нужно хранить как обычные реляционные данные, так и векторные представления фрагментов. Использование одной СУБД упрощает связь между документами, фрагментами и результатами поиска. Векторное расстояние вычисляется рядом с данными корпуса, а фильтрация по версии и типу корпуса выполняется в тех же SQL-запросах. Для этого проекта такой вариант практичнее, чем разделять реляционное хранилище и отдельное векторное хранилище, потому что источники и история должны оставаться связанными.

Библиотека sentence-transformers используется для локального построения векторных представлений моделью \texttt{BAAI/bge-m3}. В рабочем режиме приложение не отправляет корпус документации во внешний сервис для векторизации. Векторные представления рассчитываются локально при переиндексации и сохраняются в базе данных. Groq API используется только на этапе формирования ответа по найденному контексту. Вызов внешней модели вынесен в отдельный адаптер генерации, поэтому поиск, фильтрация версии и проверка подтверждений остаются в локальной части приложения.

\section{Структура программного проекта}

Структура программной части разделена по слоям. Сокращённое представление приведено в таблице~\ref{tab:impl-structure}. В таблицу включены только основные зоны ответственности, поскольку подробное перечисление внутренних файлов не несёт самостоятельной инженерной ценности для отчёта.

\begingroup
\centering
\small
\setstretch{1}\selectfont
\setlength{\tabcolsep}{4pt}
\begin{longtable}{|P{0.30\textwidth}|P{0.58\textwidth}|}
\caption{Структура программного проекта}\label{tab:impl-structure}\\
\hline
\tableheadcell{Путь} & \tableheadcell{Назначение} \\ \hline
\endfirsthead
\multicolumn{2}{l}{\tablepartlabel{продолжение таблицы \thetable}}\\
\hline
\endhead
\endfoot
\texttt{app/main.py} & Создание FastAPI-приложения, подключение маршрутов, статических файлов, шаблонов и начального заполнения версий. \\ \hline
\texttt{app/api/routes} & Маршруты пользовательского запроса, версий, истории, диагностики и административной переиндексации. \\ \hline
\texttt{app/schemas} & Pydantic-схемы запросов, ответов, истории, версий, диагностики и административных операций. \\ \hline
\texttt{app/db} & ORM-модели, перечисления и создание сессии базы данных. \\ \hline
\texttt{app/repositories} & SQL-запросы для поиска фрагментов, индексации, истории и версий. \\ \hline
\texttt{app/services} & Обработка пользовательского запроса, анализ вопроса, поиск фрагментов, повторное ранжирование, история и администрирование. \\ \hline
\makecell[l]{\texttt{app/services/}\\\texttt{ingestion}} & Загрузка, очистка, разбиение и индексация официального и дополнительного корпусов. \\ \hline
\makecell[l]{\texttt{app/services/}\\\texttt{adapters}} & Адаптеры векторных представлений и генерации через Groq API. \\ \hline
\texttt{app/web} & Шаблоны HTML, стили и JavaScript пользовательского интерфейса. \\ \hline
\texttt{tests} & Автоматические модульные, API, сценарные и интеграционные тесты. \\ \hline
\end{longtable}
\par\endgroup

\section{Реализация модели базы данных}

Модель базы данных реализована в \texttt{app/db/models.py}. В ней определены таблицы \texttt{versions}, \texttt{documents}, \texttt{chunks}, \texttt{embeddings}, \texttt{query\_history} и \texttt{query\_sources}. Набор допустимых технических значений вынесен в перечисления для типов корпуса, режимов ответа и статусов обработки.

Основной фрагмент ORM-моделей приведён в листинге~\ref{listing:impl-db-models}. В листинге оставлены только поля, которые показывают версионную привязку, фрагменты, векторы и историю запросов.

\begin{listingblock}
\captionof{listing}{Фрагмент ORM-моделей базы данных}
\label{listing:impl-db-models}
\begin{minted}{python}
class Version(Base):
    __tablename__ = "versions"
    id = mapped_column(UUID(as_uuid=True), primary_key=True, default=uuid.uuid4)
    major_version = mapped_column(String(10), unique=True, nullable=False)
    docs_base_url = mapped_column(String(255), nullable=False)
    is_supported = mapped_column(Boolean, nullable=False, default=True)

class Document(Base):
    __tablename__ = "documents"
    version_id = mapped_column(UUID(as_uuid=True), ForeignKey("versions.id", ondelete="CASCADE"), nullable=False)
    title = mapped_column(String(500), nullable=False)
    source_url = mapped_column(String(700), nullable=False)
    corpus_type = mapped_column(String(30), nullable=False)

class Chunk(Base):
    __tablename__ = "chunks"
    document_id = mapped_column(UUID(as_uuid=True), ForeignKey("documents.id", ondelete="CASCADE"), nullable=False)
    section_path = mapped_column(String(700), nullable=False)
    chunk_text = mapped_column(Text, nullable=False)
\end{minted}
\end{listingblock}

\begin{listingblock}
\listingpartlabel{продолжение листинга~3.1}
\begin{minted}[firstnumber=18]{python}
    pedagogical_role = mapped_column(String(30), nullable=False, default="overview")
class Embedding(Base):
    __tablename__ = "embeddings"
    chunk_id = mapped_column(UUID(as_uuid=True), ForeignKey("chunks.id", ondelete="CASCADE"), nullable=False, unique=True)
    embedding_model = mapped_column(String(120), nullable=False)
    embedding_dimension = mapped_column(Integer, nullable=False)
    embedding = mapped_column(Vector(settings.embedding_dimension), nullable=False)
\end{minted}
\end{listingblock}

История запросов хранится отдельно от корпуса. В таблице \texttt{query\_history} сохраняются исходный вопрос, выбранная версия, режим, текстовый ответ, JSON-структура обучающего результата, статус и задержка, а таблица \texttt{query\_sources} хранит использованные источники: позицию, оценку, роль, заголовок, URL, путь раздела и тип корпуса.

\section{Реализация миграций и физической схемы PostgreSQL}

Физическая схема создаётся миграциями Alembic. Начальная миграция создаёт расширение \texttt{vector}, таблицы корпуса и истории, ограничения допустимых значений, индексы и векторный индекс \texttt{ivfflat}. Миграция от 4 мая 2026 года переводит векторное поле на \texttt{vector(1024)} для модели \texttt{BAAI/bge-m3}. Миграции от 5 и 7 мая 2026 года фиксируют актуальную схему режима истории: три действующих режима ответа и отсутствие устаревшего поля \texttt{extended\_mode}.

Ключевые строки миграции показаны в листинге~\ref{listing:impl-alembic}. Они подтверждают использование PostgreSQL, pgvector, JSONB и индекса для векторного поиска.

\begin{listingblock}
\captionof{listing}{Ключевые элементы миграции PostgreSQL}
\label{listing:impl-alembic}
\begin{minted}{python}
embedding_dim = int(os.getenv("EMBEDDING_DIMENSION", "1024"))
op.execute("CREATE EXTENSION IF NOT EXISTS vector")

op.create_table(
    "embeddings",
    sa.Column("chunk_id", postgresql.UUID(as_uuid=True), nullable=False, unique=True),
    sa.Column("embedding_model", sa.String(length=120), nullable=False),
    sa.Column("embedding_dimension", sa.Integer(), nullable=False),
    sa.Column("embedding", Vector(embedding_dim), nullable=False),
)

op.create_table(
    "query_history",
    sa.Column("mode", sa.String(length=30), nullable=False),
    sa.Column("tutorial_json", postgresql.JSONB(astext_type=sa.Text()), nullable=True),
)

op.execute(
    "CREATE INDEX IF NOT EXISTS ix_embeddings_embedding ON embeddings "
    "USING ivfflat (embedding vector_cosine_ops) WITH (lists = 100)"
)
\end{minted}
\end{listingblock}

В Docker Compose используется образ \texttt{pgvector/pgvector:pg16}: контейнерная база данных работает на PostgreSQL 16 с установленным расширением pgvector. Версии PostgreSQL 16, 17 и 18 относятся к документации, по которой строится корпус, а не к версии серверного контейнера БД.

\section{Реализация маршрутов программного интерфейса}

Маршруты API подключаются через общий маршрутизатор с префиксом \texttt{/api}. Реализация содержит маршруты \texttt{/api/ask}, \texttt{/api/versions}, \texttt{/api/history}, \texttt{/api/health}, \texttt{/api/health/embeddings} и \texttt{/api/admin/reindex}. Дополнительно в \texttt{app/main.py} реализованы HTML-страницы \texttt{/} и \texttt{/history}, а также диагностические маршруты \texttt{/health} и \texttt{/health/embeddings} без префикса \texttt{/api}.

Контракт запроса \texttt{/api/ask} задаётся схемой \texttt{AskRequest}. Фрагмент схемы приведён в листинге~\ref{listing:impl-ask-schema}. Он показывает реальные режимы ответа и обязательную валидацию версии.

\begin{listingblock}
\captionof{listing}{Фрагмент схемы пользовательского запроса}
\label{listing:impl-ask-schema}
\begin{minted}{python}
class AskRequest(BaseModel):
    model_config = ConfigDict(extra="ignore")
    question: str = Field(min_length=3, max_length=4000)
    pg_version: str = Field(min_length=1, max_length=10)
    answer_mode: Literal["short", "detailed", "tutorial"]
    @field_validator("pg_version")
    @classmethod
    def validate_version(cls, value: str) -> str:
        normalized = value.strip()
        if not normalized.isdigit():
            raise ValueError("pg_version must be a major numeric version, e.g. '16'")
        if normalized not in settings.supported_versions:
            raise ValueError("pg_version must be one of supported versions")
        return normalized
\end{minted}
\end{listingblock}

Схема ответа \texttt{AskResponse} содержит текстовый ответ, обучающую структуру и список источников. Для текстовых режимов используется поле ответа, для обучающего режима -- вложенная обучающая структура. Примеры запросов и ответов программного интерфейса приведены в листингах~\ref{listing:appendix-request-short},~\ref{listing:appendix-request-detailed},~\ref{listing:appendix-request-tutorial},~\ref{listing:appendix-response-short} и~\ref{listing:appendix-response-tutorial}.

\begin{listingblock}[25mm]
\captionof{listing}{Пример запроса в кратком режиме}
\label{listing:appendix-request-short}
\begin{minted}[breaklines=true,breakanywhere=true,fontsize=\small]{text}
POST /api/ask
Content-Type: application/json
\end{minted}
\end{listingblock}

\begin{listingblock}[24mm]
\listingpartlabel{продолжение листинга~3.4}
\begin{minted}[breaklines=true,breakanywhere=true,fontsize=\small,firstnumber=3]{text}
{
  "question": "Как включить логическую репликацию?",
  "pg_version": "16",
  "answer_mode": "short"
}
\end{minted}
\end{listingblock}

\begin{listingblock}[65mm]
\captionof{listing}{Пример запроса в подробном режиме}
\label{listing:appendix-request-detailed}
\begin{minted}[breaklines=true,breakanywhere=true,fontsize=\small]{text}
POST /api/ask
Content-Type: application/json

{
  "question": "Подробно объясни VACUUM в PostgreSQL",
  "pg_version": "16",
  "answer_mode": "detailed"
}
\end{minted}
\end{listingblock}

\begin{listingblock}[65mm]
\captionof{listing}{Пример запроса в обучающем режиме}
\label{listing:appendix-request-tutorial}
\begin{minted}[breaklines=true,breakanywhere=true,fontsize=\small]{text}
POST /api/ask
Content-Type: application/json

{
  "question": "Объясни EXPLAIN ANALYZE простыми словами",
  "pg_version": "16",
  "answer_mode": "tutorial"
}
\end{minted}
\end{listingblock}

\begin{listingblock}[25mm]
\captionof{listing}{Пример ответа в кратком режиме}
\label{listing:appendix-response-short}
\begin{minted}[breaklines=true,breakanywhere=true,fontsize=\small]{json}
{
  "answer_mode": "short",
  "pg_version": "16",
  "answer": "...",
  "sources": [
\end{minted}
\end{listingblock}

\begin{listingblock}[24mm]
\listingpartlabel{продолжение листинга~3.7}
\begin{minted}[breaklines=true,breakanywhere=true,fontsize=\small,firstnumber=6]{json}
    {
      "title": "PostgreSQL 16 Documentation",
      "url": "https://www.postgresql.org/docs/16/...",
      "corpus_type": "official",
      "source_role": "base",
      "section_path": "Chapter 31 / Logical Replication",
      "rank_position": 1,
      "similarity_score": 0.91234
    }
  ]
}
\end{minted}
\end{listingblock}

\begin{listingblock}[24mm]
\captionof{listing}{Пример ответа в обучающем режиме}
\label{listing:appendix-response-tutorial}
\begin{minted}[breaklines=true,breakanywhere=true,fontsize=\small]{json}
{
  "answer_mode": "tutorial",
  "pg_version": "16",
  "tutorial": {
    "short_explanation": "...",
    "prerequisites": ["..."],
    "steps": ["..."],
    "notes": ["..."]
  },
  "sources": [
    {
      "title": "PostgreSQL 16 Documentation",
      "url": "https://www.postgresql.org/docs/16/...",
      "corpus_type": "official",
      "source_role": "base",
      "section_path": "Chapter 31 / Logical Replication",
      "rank_position": 1,
      "similarity_score": 0.94211
    },
    {
      "title": "Tutorial note",
      "url": "supplementary://curated/16/guide.md",
      "corpus_type": "supplementary",
      "source_role": "supplementary",
      "section_path": "Guide / Logical Replication",
      "rank_position": 2,
\end{minted}
\end{listingblock}

\begin{listingblock}[35mm]
\listingpartlabel{продолжение листинга~3.8}
\begin{minted}[breaklines=true,breakanywhere=true,fontsize=\small,firstnumber=27]{json}
      "similarity_score": 0.72111
    }
  ]
}
\end{minted}
\end{listingblock}

\section{Реализация подготовки и индексации корпуса}

Официальный корпус загружается из локальной папки \texttt{corpus/postgres/html/<version>}. Загрузчик выбирает файлы \texttt{.html} и \texttt{.htm}, строит канонический \texttt{source\_url} на основе \texttt{OFFICIAL\_DOCS\_BASE\_URL} и выбранной версии, а при ограничении \texttt{max\_pages} применяет приоритет для страниц логической репликации, конфигурации, SQL-команд и примечаний к выпускам.

Индексировать исходный HTML без очистки нельзя, потому что страницы документации содержат не только основной текст. В них присутствуют навигационные элементы, ссылки на соседние страницы, формы поиска, служебные блоки, заголовки сайта, нижние колонтитулы, сценарии и стили. Если такие элементы попадут в корпус, они будут участвовать в построении векторных представлений и могут исказить поиск. Например, повторяющаяся навигация встречается на многих страницах и создаёт одинаковый шум в разных фрагментах. Поэтому парсер сначала удаляет служебные элементы, а затем извлекает содержательные блоки документации.

При обработке HTML система работает не с плоским текстом, а со структурой страницы. Заголовки разных уровней формируют путь раздела, который сохраняется в поле \texttt{section\_path}. Обычные абзацы, элементы списков, таблицы и блоки кода преобразуются в отдельные текстовые блоки с указанием типа содержимого. Это нужно не только для отображения источников, но и для качества поиска. Фрагмент с путём раздела легче связать с темой вопроса, чем фрагмент без контекста. Например, одинаковый термин может встречаться в описании SQL-команды, параметра конфигурации или системного представления, но путь раздела помогает различить эти случаи.

Разбиение на фрагменты выполняется после выделения структурных блоков. Система не объединяет разные пути разделов и разные типы содержимого в один фрагмент. Это уменьшает смысловое смешение: абзац с описанием команды, таблица параметров и блок SQL-кода могут относиться к одной странице, но выполнять разные функции. Ограничение \texttt{CHUNK\_SIZE} задаёт максимальный размер фрагмента, а \texttt{CHUNK\_OVERLAP} переносит часть предыдущего текста в следующий фрагмент при переполнении. Перекрытие нужно для случаев, когда смысловое объяснение находится на границе двух фрагментов. Без него часть контекста могла бы потеряться при поиске и передаче данных в генеративную модель.

Для каждого фрагмента дополнительно определяется педагогическая роль. В реализации используются роли \texttt{overview}, \texttt{prerequisite}, \texttt{step}, \texttt{example} и \texttt{warning}. Они определяются по маркерам в тексте и затем учитываются при повторном ранжировании. В обучающем режиме фрагменты с примерами, шагами и предварительными условиями могут быть полезнее для формирования пошагового ответа. В кратком и подробном режимах больший вес обычно получают обзорные и предупреждающие фрагменты, потому что они чаще содержат определения, ограничения и правила использования.

Дополнительный корпус обрабатывается отдельно от официального. Для файлов из каталога \texttt{curated/<version>} проверяется поле \texttt{pg\_version} в метаданных. Если документ относится к другой версии, он не попадает в индекс выбранной версии. Это правило нужно по той же причине, что и фильтрация официальных URL: учебный материал не должен незаметно смешивать сведения разных веток PostgreSQL. Обработанные Markdown-файлы из \texttt{processed\_html} индексируются только при наличии признака \texttt{indexable}. Служебные отчёты, манифесты, резервные копии и файлы кэша исключаются, потому что они описывают состояние корпуса или процесс его подготовки, но не являются пользовательскими источниками знаний по PostgreSQL.

Индексация выполняется классом \texttt{IndexingPipeline}. Он получает или создаёт запись версии, подготавливает документы, вычисляет векторные представления пакетами, удаляет старый корпус выбранного типа и сохраняет новые документы, фрагменты и векторы. В рабочем режиме используется \texttt{EMBEDDING\_PROVIDER=local}, \texttt{EMBEDDING\_MODEL=BAAI/bge-m3}, \texttt{EMBEDDING\_DIMENSION=1024}, \texttt{EMBEDDING\_MAX\_SEQ\_LENGTH=8192}.

\section{Реализация обработки пользовательского запроса}

Обработка запроса реализована в \texttt{AskOrchestrationService}. Сервис получает объект \texttt{AskRequest}, проверяет существование версии, анализирует вопрос, выбирает корпуса, получает векторное представление вопроса, запускает поиск фрагментов, повторное ранжирование, проверку подтверждений и генерацию. После обработки сервис сохраняет историю запроса и использованные источники.

Анализ вопроса реализован в \texttt{app/services/query\_processing.py}. На этом этапе текст нормализуется: приводится к единому регистру, очищается от лишних пробелов, используется единое представление буквы <<е>>. После этого определяется намерение вопроса. В реализации используются значения \texttt{factual}, \texttt{definition}, \texttt{how\_to}, \texttt{configuration}, \texttt{comparison}, \texttt{support\_check} и \texttt{release\_changes}. Эти значения не показываются пользователю, но влияют на поиск и повторное ранжирование. Например, вопрос о поддержке или совместимости требует более строгих источников, чем общий вопрос с определением.

Отдельно выделяются технические термины и их варианты. Вопрос пользователя может быть сформулирован на русском языке, а документация PostgreSQL часто содержит английские названия команд, параметров и системных представлений. Поэтому для ряда тем заданы варианты терминов: \texttt{pg\_stat\_activity}, \texttt{EXPLAIN}, \texttt{VACUUM}, \texttt{ANALYZE}, \texttt{logical replication}, \texttt{publication}, \texttt{subscription}, \texttt{wal\_level} и другие. Эти термины используются дважды: как подсказки при построении векторного представления вопроса и как набор выражений для лексической ветки поиска.

Для некоторых тем система определяет тематический профиль. Такой профиль задаёт не только список терминов, но и положительные и отрицательные якоря. Например, для вопросов про активные запросы важны \texttt{pg\_stat\_activity}, \texttt{query\_start}, \texttt{state} и \texttt{pid}. Для вопросов про логическую репликацию важны \texttt{logical replication}, \texttt{publication}, \texttt{subscription}, параметры репликации и термины издателя/подписчика (publisher/subscriber). Если найденный фрагмент содержит похожие слова, но относится к другой теме, повторное ранжирование должно понизить его оценку. Это особенно важно для технической документации, где одинаковые слова могут встречаться в разных разделах.

Выбор корпусов зафиксирован статическим методом \texttt{resolve\_corpora}. Для краткого и подробного режимов возвращается только официальный корпус (\texttt{official}); для обучающего режима возвращаются официальный и дополнительный корпуса (\texttt{official}, \texttt{supplementary}). Если обучающий поиск не вернул дополнительных кандидатов, сервис отдельно пытается получить их из дополнительного корпуса, но только после основного поиска официального корпуса.

\section{Реализация поиска и учёта версии}

Поиск фрагментов реализован в \texttt{RetrievalRepository} и \texttt{RetrievalService}. Репозиторий извлекает кандидаты двумя ветками: векторной и лексической. Обе ветки фильтруют документы по \texttt{Version.major\_version} и списку допустимых типов корпуса. Векторная ветка сортирует результаты по \texttt{cosine\_distance}; лексическая ветка учитывает совпадения в \texttt{title}, \texttt{section\_path} и \texttt{chunk\_text}.

Векторная ветка поиска отвечает за смысловую близость вопроса и фрагментов. Пользователь может не знать точного названия раздела или команды, но описать задачу своими словами. В этом случае векторное представление вопроса сравнивается с векторными представлениями фрагментов, сохраненными в таблице \texttt{embeddings}. В SQL-запросе используется косинусное расстояние, а результаты сортируются от наиболее близких к менее близким. При этом поиск сразу ограничивается выбранной версией PostgreSQL и допустимым типом корпуса.

Лексическая ветка нужна для точных технических совпадений. В документации PostgreSQL многие важные сведения выражены через конкретные идентификаторы: имена SQL-команд, системных представлений, параметров конфигурации, функций и утилит. Векторный поиск может найти близкий по смыслу фрагмент, но для инженерного ответа важно не потерять точные совпадения вроде \texttt{work\_mem}, \texttt{pg\_stat\_activity}, \texttt{wal\_level} или \texttt{ALTER SUBSCRIPTION}. Поэтому лексическая ветка ищет расширенные термины в заголовке документа, адресе источника, пути раздела и тексте фрагмента. Совпадения в заголовке и пути раздела получают больший вес, чем совпадения только в основном тексте.

Результаты двух веток объединяются по идентификатору фрагмента. Если один и тот же фрагмент найден и семантически, и лексически, он не должен дублироваться в списке кандидатов. После объединения применяется предварительная оценка, которая сочетает семантическое сходство и лексический сигнал. Эта оценка не используется как окончательный вывод о качестве ответа. Она нужна для сокращения набора кандидатов перед более подробным повторным ранжированием.

Формула предварительной оценки показана в листинге~\ref{listing:impl-prerank}. Она используется только для сортировки кандидатов перед повторным ранжированием.

\begin{listingblock}
\captionof{listing}{Формула предварительного ранжирования кандидатов}
\label{listing:impl-prerank}
\begin{minted}{python}
semantic_similarity = max(0.0, 1.0 - (item.distance / 2.0))
lexical_signal = min(item.lexical_score / 8.0, 1.0)
pre_rank = 0.74 * semantic_similarity + 0.26 * lexical_signal
\end{minted}
\end{listingblock}

Контроль версии реализован методом \texttt{\_enforce\_version\_guard} класса \texttt{RetrievalService}. Сервис удаляет официальный источник, если в его URL нет ожидаемого фрагмента \texttt{/docs/<pg\_version>/}. Тем самым SQL-фильтр по версии дополняется явной проверкой адреса официальной документации. Фрагменты реализации контроля версии и выбора корпусов приведены в листингах~\ref{listing:appendix-version-guard} и~\ref{listing:appendix-corpora}.

\begin{listingblock}[24mm]
\captionof{listing}{Реализация контроля версии источника}
\label{listing:appendix-version-guard}
\begin{minted}{python}
@staticmethod
def _enforce_version_guard(
    *, rows: list[RetrievedChunk], pg_version: str
) -> list[RetrievedChunk]:
    expected_token = f"/docs/{pg_version}/"
    filtered: list[RetrievedChunk] = []

    for item in rows:
        if (
            item.corpus_type == CorpusType.OFFICIAL.value
            and expected_token not in item.source_url
        ):
            continue
        filtered.append(item)
    return filtered
\end{minted}
\end{listingblock}

\enlargethispage{18mm}
\begin{listingblock}[24mm]
\captionof{listing}{Правило выбора корпусов}
\label{listing:appendix-corpora}
\begin{minted}{python}
@staticmethod
def resolve_corpora(answer_mode: str) -> list[str]:
    if answer_mode == ModeType.TUTORIAL.value:
\end{minted}
\end{listingblock}

\begin{listingblock}[24mm]
\listingpartlabel{продолжение листинга~3.11}
\begin{minted}[firstnumber=4]{python}
        return [CorpusType.OFFICIAL.value, CorpusType.SUPPLEMENTARY.value]
    return [CorpusType.OFFICIAL.value]
\end{minted}
\end{listingblock}

\section{Реализация повторного ранжирования и проверки подтверждений}

Повторное ранжирование реализовано в \texttt{RerankingService}. Для каждого кандидата вычисляется итоговая оценка на основе семантического сходства, лексического пересечения, совпадения с заголовком и путём раздела, совпадения технических терминов, лексического сигнала, типа корпуса и педагогической роли. Для логической репликации, вопросов совместимости и параметров конфигурации применяются дополнительные бонусы и штрафы.

Повторное ранжирование нужно потому, что первичный поиск возвращает кандидатов, но не всегда в порядке, удобном для ответа. Векторная близость может поднять общий тематический фрагмент, а лексическое совпадение может поднять фрагмент, где термин встречается случайно. Поэтому итоговая оценка строится из нескольких признаков. Семантическое сходство показывает общую близость вопроса и текста. Лексическое пересечение учитывает общие токены вопроса и фрагмента. Совпадение с заголовком и путём раздела помогает поднять страницы, где тема выражена в структуре документации. Совпадение технических терминов проверяет, что найденный фрагмент содержит именно те идентификаторы и понятия, которые важны для вопроса.

Тип корпуса участвует в оценке отдельно. Официальный корпус получает положительный приоритет, а дополнительный корпус ограничивается. Это особенно важно для обучающего режима: учебные материалы могут помочь в подаче, но не должны вытеснять официальные разделы документации. Также сервис ограничивает число фрагментов из одного документа. Если этого не делать, контекст может состоять из нескольких соседних фрагментов одной страницы, а другие полезные источники не попадут в результат. Ограничение уменьшает повторы и делает список источников более разнообразным.

Для отдельных тем применяются дополнительные правила. Вопросы о логической репликации требуют якорей, связанных с публикациями, подписками, слотами репликации и параметрами настройки. Вопросы о поддержке или совместимости требуют источников, где есть признаки ограничений, версий, издателя/подписчика или основной версии. Вопросы об изменениях версии должны опираться на примечания к выпускам, совместимость или миграционные сведения, а не на произвольные страницы справочника команд. Эти правила не заменяют поиск, а корректируют порядок кандидатов после первичного отбора.

В кратком и подробном режимах итоговый список строится только из официальных фрагментов. В обучающем режиме дополнительный корпус может быть добавлен к официальным источникам, но не заменяет их. Практическая реализация ограничивает число фрагментов одного документа, чтобы контекст не состоял из повторов одной страницы.

Проверка достаточности подтверждений реализована методом \texttt{\_has\_sufficient\_evidence} и выполняется до вызова Groq API. Сервис проверяет наличие официальных источников, пороги итоговой оценки, семантическое сходство, лексическое пересечение, совпадение заголовка и технических терминов. Для логической репликации проверка строже: дополнительно требуются тематические якоря, потому что в этой теме легко спутать логическую репликацию с физической репликацией, резервными серверами, \texttt{LISTEN/NOTIFY} или утилитами работы с WAL. Если найденный контекст связан с похожими словами, но не содержит нужных подтверждений, система не передаёт его модели как основание для уверенного ответа.

Если проверка не пройдена, приложение возвращает безопасный отказ. Для краткого и подробного режимов это текст о том, что в найденных фрагментах документации выбранной версии недостаточно тематического подтверждения. Для обучающего режима возвращается структура с кратким пояснением, пустыми шагами и замечаниями о необходимости уточнить вопрос. Такой результат лучше, чем связный, но неподтверждённый ответ, потому что пользователь сразу видит ограничение найденного контекста.

После успешной проверки сервис формирует контекст для генеративной модели. В него включаются не все найденные фрагменты, а отобранные источники с заголовком, путём раздела, причиной релевантности и выдержкой текста. Системная инструкция требует отвечать на русском языке, не смешивать версии PostgreSQL, не выводить служебные поля и не добавлять источники в текст ответа, потому что источники возвращаются отдельной структурой API. После получения ответа текст дополнительно очищается от лишних ссылок, служебных блоков и повторов. Затем результат сохраняется в истории вместе с источниками и задержкой обработки.

\section{Реализация пользовательского интерфейса}

Пользовательский интерфейс создан шаблонами \texttt{app/web/templates/index.html} и \texttt{app/web/templates/history.html}, стилями \texttt{app/web/static/style.css} и сценариями \texttt{app/web/static/app.js}, \texttt{app/web/static/history.js}. Главная страница содержит поле вопроса, быстрые примеры, выбор версии PostgreSQL, сегмент выбора режима ответа, блок результата, блок источников и список недавних запросов.

Структура главной страницы подчинена основному сценарию: пользователь формулирует вопрос, выбирает версию PostgreSQL и режим ответа, после чего получает результат в той же рабочей области. Интерфейс не содержит отдельного переключателя корпуса. Это согласуется с серверной логикой: состав корпуса определяется режимом ответа, а не произвольным выбором пользователя. Такой вариант снижает риск ошибки, когда пользователь случайно включает дополнительный корпус для краткого ответа и получает менее проверяемый результат.

JavaScript главной страницы загружает список версий через \texttt{/api/versions}, формирует структуру запроса для \texttt{/api/ask}, обрабатывает ответы трёх режимов и отображает источники. Для краткого и подробного режимов используется текстовый блок ответа. Для обучающего режима отображаются отдельные секции: краткое объяснение, предварительные условия, пошаговые действия, замечания и частые ошибки. Пустые секции скрываются, поэтому пользователь не видит технических заготовок, если сервер вернул неполную обучающую структуру.

Блок загрузки показывает смену этапов обработки: поиск источников, формирование ответа и проверку формулировки. Это не добавляет новых функций, но помогает пользователю понимать состояние длительного запроса. При ошибке интерфейс выводит нейтральное сообщение и сохраняет технические детали в раскрываемом блоке. Такой подход отделяет пользовательское объяснение от диагностической информации: обычный пользователь видит причину сбоя в краткой форме, а разработчик может скопировать детали для отладки.

Источники отображаются отдельными карточками. Для каждого источника показываются название, путь раздела, тип корпуса, версия, ссылка и метаданные релевантности. Если источников много, интерфейс сначала показывает ограниченное число карточек и дает возможность раскрыть остальные. Благодаря этому ответ остается в фокусе, но трассировка до документов не скрывается.

Страница истории обращается к \texttt{/api/history?limit=100} и показывает вопрос, дату, версию, режим, статус, краткий предварительный просмотр ответа и раскрываемый список источников. На главной странице боковая панель истории запрашивает только четыре последних запроса через \texttt{/api/history?limit=4}. Поэтому пользователь может быстро вернуться к недавнему вопросу, а полный просмотр истории вынесен на отдельную страницу.

Руководство пользователя по работе с интерфейсом приведено в приложении~Г.

\section{Выводы}

В третьей главе описана реализация веб-приложения: выбранный набор технологий, структура проекта, ORM-модели, миграции PostgreSQL, маршруты программного интерфейса, подготовка и индексация корпуса, обработка пользовательского запроса, поиск фрагментов с учётом версии, повторное ранжирование, проверка достаточности подтверждений и пользовательский интерфейс. Раздел тестирования вынесен в отдельную четвёртую главу, поэтому вывод по главе относится только к программной части.
